\begin{proposition}[Finite rank drop of a normalization]\hypertarget{tag-02WM}{}\label{night:P4.15}\label{stmt:P4.15}
\StatementRecord{night:P4.15}{proved}{02WM}{\SeriesRef{cdp:st:A.5.3}, \SeriesRef{night:P4.12}, \SeriesRef{night:P4.13}}

Let $f:Y\to B$ be a proper surjective holomorphic map with
connected fibres from a connected smooth compact complex threefold
with $b_1(Y)=b_2(Y)=0$ to a smooth compact curve, and suppose that
its smooth general fibre is a complex two-torus. Let $mS$ be a full
fibre whose reduced support $S$ is irreducible and singular.
Suppose its normalization $\nu:T\to S$ has smooth domain $T$. Set
$j:T\xrightarrow{\nu}S\hookrightarrow Y$ and assume that
$dj:T_T\to j^*T_Y$ has rank two away from finitely many points.
Let $\ell$ be the length of its maximal-minor scheme. Then
\[
 e(T)=K_T^2-\ell.
\]
The alternatives for $T$ are a rational surface with
$e(T)=6-\ell/2$ and $0\leq\ell\leq6$, a minimal ruled surface
over an elliptic curve with $\ell=0$, or a Hopf surface with
$\ell=0$. In either of the last two alternatives $e(S)=0$.
In those last two alternatives, the reduced conductor $R\subset T$
and $C=\operatorname{Sing}(S)_{\mathrm{red}}$ are disjoint unions of
smooth elliptic curves, and the restriction $\nu|_R:R\to C$ is
finite and unramified.
\end{proposition}
\begin{proof}

\begin{enumerate}[label={$\langle1\rangle$\arabic*.},leftmargin=*]
\item The differential is injective as a sheaf map: its kernel vanishes
on a dense open set, and its locally free source is torsion-free.
Its rank-one cokernel $Q$ is locally free away from finitely many
points. The depth lemma applied to the length-one locally free
resolution gives depth at least one at those points. Hence $Q$
has no finite-length subsheaf, and all its possible torsion would
be supported there. Thus $Q$ is torsion-free, and $L=Q^{**}$ is a
line bundle. In local frames the signed maximal-minor row of $dj$
annihilates its two columns. Its determinant transition functions
give a map to $L=\det(j^*T_Y)\otimes\det(T_T)^{-1}$.
The induced map $Q\to L$ is generically injective, hence injective
because $Q$ is torsion-free. Its image is the ideal generated by
those minors, tensored with $L$. We obtain the exact sequence
\[
 0\longrightarrow T_T\xrightarrow{dj}j^*T_Y
   \longrightarrow\mathcal I_Z\otimes L\longrightarrow0,
 \qquad L=K_T\otimes j^*(\det T_Y),\qquad
 \operatorname{length}(Z)=\ell.
\]
Poincar\'e duality gives $b_4(Y)=b_2(Y)=0$, so
$c_1(T_Y)=c_2(T_Y)=0$ in rational cohomology. Writing $k=c_1(K_T)$,
the Chern-character identity for this sequence is
\[
 3=\left(2-k+\frac{k^2-2c_2(T_T)}2\right)
     +\left(1+k+\frac{k^2}2-[Z]\right).
\]
Its degree-four part gives $e(T)=K_T^2-\ell$, and Noether's formula
gives $2e(T)+\ell=12\chi(\mathcal O_T)$.

Let $D$ be the full conductor divisor on $T$, and put $R=D_{\mathrm{red}}$.
Propositions~\ref{cdp:st:A.5.3} and~\ref{night:P4.12}, with the
identity as the resolution of $T$, give
$\omega_T(D)\simeq\nu^*\omega_S$ and $\omega_S^{\otimes m}\simeq\mathcal O_S$.
In particular $K_T+D\equiv0$. Proposition~\ref{night:P4.12}
also gives $\kappa(T)=-\infty$ and nonnormality of the singular $S$.
Since the hypersurface $S$ satisfies $S_2$, its nonnormal locus
contains a curve; thus $D$ is nonempty.
Let $q:T\to T_0$ contract $r$ point blowups to a minimal model.
Pluricanonical sections descend across each blowup by Hartogs'
extension on the complement of its centre, so $\kappa(T_0)=-\infty$
and $p_g(T_0)=0$.

For a K\"ahler $T_0$, the equality $p_g=0$ makes every real
degree-two cohomology class of type $(1,1)$. A sufficiently close
rational perturbation of a K\"ahler class, followed by integral
scaling, gives an integral K\"ahler class. Kodaira's embedding
theorem \cite[Main Theorem, p.~314]{Kodaira1954} therefore gives an
embedding $\Phi:T_0\hookrightarrow\mathbb P^N$. Apply
\cite[VI.17--VI.18, pp.~111--112; V.10, pp.~79--80; III.10, p.~40]{Beauville1978}
to $\Phi(T_0)$ and pull its ruling back by $\Phi$: the model is
rational or a projective-line bundle over a positive-genus curve.

For a non-K\"ahler $T_0$, \cite[Theorem~5.1, pp.~27--29]{OVV2020}
gives an isotrivial elliptic surface or a class-VII surface, with
overlap permitted. In the elliptic alternative, the local maps in
\cite[Theorem~4.2 and Remark~4.4, pp.~17--21]{OVV2020} are
\[
 (E\times\Delta_s)/\langle(z,s)\mapsto(z+e,\zeta_m s)\rangle
 \longrightarrow\Delta_t,\qquad t=s^m,
\]
where $e$ has exact order $m$. The elliptic action of
\cite[Proposition~4.5, pp.~21--22]{OVV2020} gives a nowhere-zero
vertical vector field. For its ruling $p:T_0\to B_0$, the quotient
of $T_{T_0}$ by this vector field maps to $p^*T_{B_0}$ with divisor
$R_p=\sum(m_i-1)E_i$, as follows from $dt=m_i s^{m_i-1}ds$.
Consequently $K_{T_0}=p^*K_{B_0}+R_p$, so $K_{T_0}^2=0$;
the vector field gives $c_2(T_0)=0$. Noether's formula and $p_g=0$
give $h^1(\mathcal O_{T_0})=1$.
The identity $b_1=2h^1(\mathcal O)-1$ in
\cite[Theorem~2.5, pp.~7--9]{OVV2020} gives $b_1(T_0)=1$.
Thus this elliptic alternative also has the class-VII invariants
$b_1=1$ and $\kappa=-\infty$.

For a projective-line bundle over a genus-$g$ curve, followed by
$r$ point blowups, the bundle and blowup formulas
\cite[II.3, pp.~17--18; III.18, pp.~46--47]{Beauville1978} give
\[
 e(T)=4(1-g)+r,\qquad K_T^2=8(1-g)-r.
\]
The Chern identity becomes $2r+\ell=4(1-g)$. For $g\ge1$ it forces
$g=1$ and $r=\ell=0$.
For a class-VII minimal model, write $b_2(T_0)=k_0\ge0$.
The equalities $b_1=1$, $p_g=0$ and
$b_1=2h^1(\mathcal O)-1$ give $\chi(\mathcal O_{T_0})=0$.
Thus $e(T)=k_0+r$ and $K_T^2=-k_0-r$, and the Chern identity gives
$2(k_0+r)+\ell=0$. Hence $T$ is minimal with $b_2=0$ and $\ell=0$.
By \cite[Theorem~1, p.~1]{BBK2019} it has a finite \emph{\'etale}
cover $a:T'\to T$ with $T'$ Hopf or Inoue. The inverse image
$a^{-1}(R)$ is a nonempty curve. Inoue surfaces contain no curves
by \cite[Proposition~2(i), pp.~274--275; Proposition~3(i), p.~277;
paragraph following (21), p.~279]{Inoue1974}, so $T'$ is Hopf.
The composite of its universal cover
$\mathbb C^2\setminus\{0\}\to T'$ with $a$ is the universal cover
of $T$, since $\mathbb C^2\setminus\{0\}$ is simply connected.
Therefore $T$ is Hopf.

For rational $T$, $\chi(\mathcal O_T)=1$ and $e(T)\ge3$ by the
same rational minimal-model and blowup formulas. Hence
$e(T)=6-\ell/2$ and $0\le\ell\le6$.
When $\ell=0$, the maximal-minor ideal is the unit ideal, so $dj$
has rank two everywhere.

\item Consider the two latter alternatives. First let $T$ be a
minimal ruled surface over an elliptic curve $E$. Its ruling
$p:T\to E$ is a $\mathbb P^1$-bundle by
\cite[III.10, p.~40]{Beauville1978}. For a ruling fibre $F$,
the equality $K_T+D\equiv0$ gives $D\cdot F=2$.
If a horizontal component $H$ occurs with multiplicity one,
adjunction gives
\[
 2p_a(H)-2=(K_T+H)H=-(D-H)H\leq0.
\]
The map to $E$ gives geometric genus at least one. Equality
therefore holds, $H$ is smooth elliptic, and $H$ is disjoint
from every other component of $D$. There is no vertical
component because every ruling fibre meets $H$. The possibilities
are one elliptic bisection, of degree two over $E$, or two
disjoint elliptic sections.

\item The other possibility is $D=2H+\sum_j a_jF_j$, where $H$ is a
section, $F_j$ are distinct ruling fibres and $a_j>0$.
The normalization of $H$ maps isomorphically to $E$ by its
degree-one ruling map. Its map into $T$ is a section of $p$,
so its image $H$ is smooth elliptic.
The reduced conductor $R$ has Euler number equal to the number $v$
of these fibres: $e(H)=0$, $e(F_j)=2$, and each meets $H$ once.
The normalization restricts to an isomorphism
$T\setminus R\simeq S_{\mathrm{reg}}$. Hence
\[
 e(S_{\mathrm{reg}})=e(T)-e(R)=-v.
\]
Proposition~\ref{night:P4.13} forces $v=0$. Thus in every case the reduced conductor is
a disjoint union of smooth elliptic curves.

\item On a Hopf surface every irreducible curve $H$ satisfies
$H^2=K_TH=0$, so $p_a(H)=1$. A rational curve would lift from
$\mathbb P^1$ to the universal cover $\mathbb C^2\setminus\{0\}$,
where both coordinate functions must be constant. Thus $H$
is smooth elliptic. Distinct curves have intersection zero and
are disjoint, proving the same conductor description.

\item At a point of $C=\operatorname{Sing}(S)_{\mathrm{red}}$,
the finite immersive normalization presents the germ as distinct
smooth surface branches $S_i$. The reduced conductor on each branch is a single
smooth curve germ $C_i$. Any two branches intersect in a curve
in the smooth ambient threefold, and its support lies in both
$C_i$ and $C_j$. Their irreducibility forces those curve germs
to agree. Thus the image singular locus is one smooth curve
germ, and the conductor maps locally isomorphically to it.
Globally the restriction $\nu|_R:R\to C$ is finite and unramified.
For a component of degree $d$, Riemann--Hurwitz gives
$0=d(2g(C_i)-2)$, so every image component is elliptic.
Euler additivity for the normalization and its conductor gives
\[
 e(S)=e(T)-e(R)+e(C)=0.
\]

\end{enumerate}
\end{proof}